% Section 4.3
\section{Efficient KV Working-Set Management}
\label{sec:efficient-kv-working-set-management}

KVMem extends the effective context capacity of a model by dynamically
managing historical KV blocks across GPU memory, CPU memory, and SSD.  Upon a
new query, KVMem retrieves a query-dependent subset of the historical blocks
and arranges them into a compact KV working set on the GPU.  A na\"ive
implementation of this design introduces three sources of overhead.  First,
blocks removed from the GPU must be synchronously written to a lower tier,
which places stage-out on the critical path of inference.  Second, consecutive
retrievals often select overlapping working sets, so rebuilding each working
set from scratch causes redundant stage-in traffic.  Third, the blocks that
must be loaded are physically scattered in host memory and still need to be
materialized at their new compact positions before they can be consumed by
attention.

We address these overheads with three complementary optimizations:
\emph{Proactive Chunk Writeback}, \emph{Retrieval-Aware Hierarchical Reuse},
and \emph{Packed and Pipelined KV Rematerialization}.  They respectively hide
stage-out, reduce the amount of data that must be loaded, and accelerate the
transfer and reconstruction of unavoidable misses.

\subsection{Proactive Chunk Writeback}
\label{sec:proactive-chunk-writeback}

\paragraph{Synchronous eviction bottleneck.}
In a na\"ive implementation, a KV block remains GPU-only until the active
working set reaches its capacity limit.  To reuse the corresponding GPU page,
the system must synchronously gather the block, transfer it from GPU to CPU,
scatter it into host storage, optionally write it to SSD, and only then release
the GPU page.  The pressure-selection path is therefore serialized as
\begin{equation}
  T_{\mathrm{evict}} =
  T_{\mathrm{gather}} +
  T_{\mathrm{D2H}} +
  T_{\mathrm{CPU\ scatter}} +
  T_{\mathrm{SSD}},
  \label{eq:synchronous-eviction}
\end{equation}
where the SSD term is present only when the CPU tier cannot retain the block.
The model cannot continue to use the reclaimed page until this sequence has
completed.

This synchronization is unnecessary for completed KV blocks.  Transformer KV
state is append-only: once all tokens of a prefill block have been processed,
the contents of that block no longer change.  Although the block may remain in
the active GPU working set, a lower-tier copy can be created before the block
is actually selected for eviction.

\paragraph{Background writeback.}
KVMem proactively writes back completed blocks at prefill-chunk boundaries.
After prefill chunk \(n\) has completed, the system starts a packed D2H
transfer for the newly completed blocks.  CPU scatter and optional SSD
persistence proceed in background workers while the GPU computes prefill
chunk \(n+1\).  The resulting schedule is
\begin{equation}
\begin{split}
  &\text{complete chunk } n
  \;\rightarrow\; \text{start writeback of chunk } n,\\
  &\text{compute chunk } n+1
  \;\parallel\; \text{D2H, CPU admission, and SSD persistence}.
\end{split}
\label{eq:proactive-writeback-schedule}
\end{equation}
Because a prefill chunk typically provides a substantially longer compute
window than the packed D2H and host-side copy, most of the writeback work is
hidden by subsequent model computation.

KVMem explicitly tracks the ownership of every block through the following
state transitions:
\begin{equation}
  \text{GPU-only}
  \;\rightarrow\;
  \text{backing-in-flight}
  \;\rightarrow\;
  \text{clean CPU/SSD backing}
  \;\rightarrow\;
  \text{GPU page releasable}.
  \label{eq:writeback-ownership}
\end{equation}
While a block is in the backing-in-flight state, its GPU page remains the
authoritative copy and cannot be reused.  The page becomes releasable only
after the transfer and lower-tier admission have completed.  Consequently, a
later pressure selection can evict a clean block by updating ownership and
page-table metadata, without performing another D2H transfer.

The writeback pipeline uses a bounded number of pinned slabs and a bounded
worker queue.  This prevents background persistence from becoming an
unbounded host-memory copy of the KV history.  If background workers fall
behind, pressure selection waits only for the incomplete records needed for
the current eviction, providing correctness-preserving backpressure.  Partial
or still-mutable blocks are not treated as clean.

Proactive writeback does not primarily reduce the number of bytes that must
eventually be stored.  Instead, it removes those bytes from the inference
critical path.  In our controlled CPU-only 512K-context experiment, aggregate
stage-out across 12 reselections decreased from \(3.645\)~s to \(0.268\)~s,
a \(92.7\%\) reduction, turning pressure-time eviction into an almost
metadata-only operation.

\subsection{Retrieval-Aware Hierarchical Reuse}
\label{sec:retrieval-aware-hierarchical-reuse}

\paragraph{Redundant loading across reselections.}
Let \(W_{t-1}\) and \(W_t\) denote the KV working sets selected by two
consecutive queries.  Reconstructing \(W_t\) from a lower tier would transfer
\(\lvert W_t\rvert\) blocks even when the two selections overlap.  Instead,
KVMem decomposes the update into retained, incoming, and outgoing blocks:
\begin{align}
  R_t &= W_t \cap W_{t-1}, \label{eq:retained-set}\\
  L_t &= W_t \setminus W_{t-1}, \label{eq:load-set}\\
  E_t &= W_{t-1} \setminus W_t. \label{eq:evict-set}
\end{align}
Only \(L_t\), rather than the complete \(W_t\), is a candidate for lower-tier
stage-in.

\paragraph{GPU working-set delta reuse.}
KVMem preserves the physical GPU pages of blocks in \(R_t\) whenever possible,
loads only blocks in \(L_t\), and releases or demotes blocks in \(E_t\).  More
generally, if \(G_{t-1}\) is the set of blocks currently resident on the GPU,
the required stage-in set is
\begin{equation}
  L_t^{\mathrm{physical}} = W_t \setminus G_{t-1}.
  \label{eq:physical-load-set}
\end{equation}
This delta update reduces transfer volume without changing retrieval scores or
the selected block IDs.

Physical residency and positional validity are distinct.  If a retained block
occupies the same compact position in the new working set, both its K and V
pages can be reused directly.  If its compact position changes, V remains
reusable because it is position independent, whereas K must be rematerialized
for the new position.  KVMem reconstructs such K blocks from a
position-independent, immutable raw-K authority instead of repeatedly
rotating a low-precision working K.  A moved block therefore avoids a complete
KV reload, although it may still incur raw-K transfer and re-RoPE work during
materialization.

\paragraph{Retrieval-aware CPU residency.}
GPU reuse cannot serve blocks that have already left the active working set.
For a GPU miss, latency depends strongly on whether the block is found in CPU
memory or must be read from SSD.  Conventional LRU replacement is poorly
matched to KVMem because its access sequence is generated by semantic
retrieval rather than by ordinary temporal locality.  A block that has not
been touched recently may still be repeatedly relevant to future queries,
while a recently produced block may never be selected again.

KVMem therefore uses retrieval history as a residency signal.  The CPU
admission and eviction policy considers whether a block was selected in recent
selection epochs, how frequently it has been selected, and whether it belongs
to a sink, recent, or explicitly pinned region.  Repeatedly selected hot
blocks are preferentially retained in CPU memory, while semantically cold
blocks are preferentially demoted to SSD.  This policy changes only the
physical source of a selected block; it does not bias or modify semantic
ranking.

The resulting hierarchy resolves every selected block in the following order:
\begin{enumerate}
  \item a GPU hit reuses the existing page and requires no data transfer;
  \item a CPU hit requires only CPU-to-GPU transfer;
  \item an SSD miss is read into CPU memory and then transferred to the GPU.
\end{enumerate}
This organization first eliminates redundant transfers through GPU delta
reuse, and then reduces slow-tier traffic through retrieval-aware CPU
residency.  In a trace containing 171 reselections, the heat-aware policy
reduced the number of SSD-loaded blocks from \(6{,}170\) to \(2{,}690\), a
\(56.4\%\) reduction in SSD traffic, and reduced aggregate stage-in from
\(6.000\)~s to \(5.058\)~s.

\subsection{Packed and Pipelined KV Rematerialization}
\label{sec:packed-pipelined-rematerialization}

\paragraph{Sparse transfer bottleneck.}
Hierarchical reuse minimizes the set of blocks that must be loaded, but the
remaining misses are physically scattered across host-memory slots, and their
destination GPU pages are also non-contiguous.  Issuing one H2D operation per
block results in many small PCIe transactions, repeated DMA setup, frequent
CUDA API calls, and poor effective link bandwidth.

Moreover, transfer completion is not the true completion point of KVMem
stage-in.  The immutable raw K stored by KVMem is independent of the final
compact position.  After selection, K must be placed in its destination page
and have RoPE applied for its new compact position before it can participate in
attention.  We therefore define rematerialization as the complete conversion
from scattered historical blocks to an executable compact KV working set.

\paragraph{Packed transfer.}
KVMem decouples logical block placement from physical PCIe transfer
granularity.  Multiple CPU workers gather scattered source blocks into a
contiguous pinned slab.  The system then performs a large contiguous H2D
transfer into a GPU staging buffer, followed by a GPU kernel that scatters the
records into their destination KV pages:
\begin{equation}
\begin{split}
  \text{scattered CPU blocks}
  &\;\rightarrow\; \text{contiguous pinned slab}\\
  &\;\rightarrow\; \text{bulk H2D}\\
  &\;\rightarrow\; \text{GPU page scatter}.
\end{split}
\label{eq:packed-stagein}
\end{equation}
CPU gather and GPU scatter add local memory operations, but these operations
are substantially cheaper than thousands of small PCIe submissions.  For SSD
misses, physically adjacent records are first coalesced into larger reads and
then enter the same packed H2D path.

\paragraph{Pipelined, position-aware reconstruction.}
Raw-K rematerialization is divided into three stages: CPU gathering, H2D
transfer, and GPU scatter with re-RoPE.  KVMem uses double-buffered pinned host
slabs and GPU staging buffers to process different batches concurrently:
\begin{equation}
\begin{split}
  &\text{CPU gather of batch } n+1\\
  &\qquad\parallel\; \text{H2D of batch } n\\
  &\qquad\parallel\; \text{GPU scatter/re-RoPE of batch } n-1.
\end{split}
\label{eq:rematerialization-pipeline}
\end{equation}
Cross-stream CUDA events protect buffer reuse without introducing a
device-wide synchronization at every batch.

The immutable raw-K authority uses a block-major layout so that the data needed
for one block is contiguous across the relevant layers, reducing random host
access during gathering.  The GPU rematerialization kernel uses a precomputed
RoPE sine/cosine table, avoiding repeated \texttt{powf} and \texttt{sincosf}
evaluation.  MTP rematerialization uses larger transfer batches to reduce
small-kernel and small-DMA overhead.  In addition, V stage-in and raw-K
materialization are scheduled concurrently when their data dependencies allow
it.

After pipeline warm-up, the exposed latency approaches
\begin{equation}
  T_{\mathrm{remat}} \approx
  \max\!\left(
    T_{\mathrm{CPU\ gather}},
    T_{\mathrm{H2D}},
    T_{\mathrm{scatter+RoPE}}
  \right)
  + T_{\mathrm{startup/drain}},
  \label{eq:pipelined-rematerialization-time}
\end{equation}
rather than the sum of the three main stages.  To preserve numerical
correctness, KVMem does not repeatedly apply positional deltas to a
low-precision working K; position changes are reconstructed from the
immutable raw-K authority.

In our controlled CPU-only 512K-context experiment, adding packed
rematerialization after proactive writeback and hierarchical reuse reduced
aggregate stage-in wall time from \(1.613\)~s to \(0.310\)~s.  Aggregate KV
assembly and re-RoPE decreased from \(28.478\)~s to \(7.901\)~s, a
\(3.60\times\) speedup across 12 reselections.

\subsection{End-to-End Working-Set Update}
\label{sec:end-to-end-working-set-update}

The three optimizations form one end-to-end working-set-management pipeline.
As prefill produces stable KV blocks, Proactive Chunk Writeback creates clean
lower-tier copies in the background.  When a new query produces \(W_t\),
Retrieval-Aware Hierarchical Reuse preserves GPU hits, resolves remaining
blocks from the CPU whenever possible, and accesses SSD only for cold misses.
Packed and Pipelined KV Rematerialization then converts those unavoidable
misses into bulk transfers and overlaps data movement with position-aware KV
reconstruction.  Finally, outgoing GPU pages with completed clean backing are
released without synchronous writeback.

The optimizations address orthogonal components of KVMem overhead:
\begin{itemize}
  \item \textbf{Proactive Chunk Writeback} hides data movement out of the GPU;
  \item \textbf{Retrieval-Aware Hierarchical Reuse} reduces the data that must
        be loaded;
  \item \textbf{Packed and Pipelined KV Rematerialization} accelerates the
        transfer and reconstruction of the remaining misses.
\end{itemize}
On the current controlled 512K-context workload, the measured KVMem management
overhead across 12 reselections decreased from \(63.957\)~s to \(8.488\)~s,
an \(86.7\%\) reduction and a \(7.53\times\) speedup.  End-to-end TTFT
decreased from \(341.844\)~s to \(287.176\)~s.  Lower-level mechanisms such as scalable one-dot
retrieval scoring, coalesced SSD I/O, persistent worker pools, and prefill-plan
reuse support the scalability of this pipeline, but are implementation details
rather than separate high-level optimizations.
